\section{Occupation of retarded bins}\label{occ:section}

We estimate how much particle mass a receiver can meet in a prescribed
momentum, angle, and distance range. Two changes of variables give crude
bounds. The signed increment assumption then limits how long a source can
remain close to the receiver. A time-averaged energy budget will make the
result summable over the small angular sector.

The residence-time and relative-velocity viewpoint is inspired by
Pfaffelmoser's Vlasov--Poisson argument \cite{Pfaffelmoser1992}.
This is methodological inspiration, not a theorem input for the
Vlasov--Maxwell system: the retarded geometry and source-acceleration
terms require the estimates proved here.

Throughout this section assume \eqref{eq:bootstrap} up to a time
$\tau\leq t_*$. Fix a supported receiver $X$, a dyad $w$, and a closed
interval $J\subset[0,\tau]$ of length at most $I$, where
\begin{equation}\label{occ:interval}
 0<I\leq(w/P)^2.
\end{equation}
We integrate over any measurable subset of
$\{t\in J:w/8\leq q_X(t)\leq8w\}$. The receiver need not remain in
this momentum range at other times of $J$.
All constants below may depend on the fixed horizon, initial energy,
$\|f_0\|_\infty$, and the fixed partition parameters, but not on $M,A,P$,
$\tau,J$, or the measurable subset. In particular, they are uniform over
supported receivers. We use the already established bounds
\[
 0\leq f\leq\|f_0\|_\infty,
 \qquad \sup_{0\leq s\leq\tau}
  \int_{\R^3\times\R^3}q(v)f(s,y,v)\,dy\,dv\leq C.
\]

\subsection{The momentum and angular ranges}

Let $H_0=\max(1,T_0)$. At a source point on the receiver's backward cone
write
\[
 r=t-s>0,\qquad y=X(t)-rn,\qquad n\in S^2,
\]
and retain the notation
\[
 u=v/q(v),\quad a=a(t),\quad
 d=1-n\cdot u,\quad D=1-n\cdot a,\quad e=1-a\cdot u.
\]
Choose smooth nonnegative dyadic partitions for
$r/H_0,\sqrt{d/2},\sqrt{D/2},\sqrt{e/2}$, with representatives
$h,\theta,\phi,\sigma\in\{1,2^{-1},2^{-2},\ldots\}$.
Each factor is supported where its argument lies between one half and
twice its representative. The partitions sum to one, including near the
upper endpoint $1$, have bounded overlap, and have uniformly bounded derivatives
with respect to the logarithm of each variable. For example, normalized dilates of a fixed bump
supported in $(1/2,2)$ and positive on $[3/4,3/2]$ give such partitions.
Use the analogous partition of $q\geq1$ with representatives
$p_0=1,2,4,\ldots$, but write $p=64p_0$.
On its support,
\[
 p/128\leq q\leq p/32,\qquad 64\leq p\leq P.
\]
The last inequality follows from the assumed bound on all supported
momenta. We use closed support ranges for upper estimates; this preserves
bounded overlap.

These are angular scales with the inverse-momentum deficits included,
not literal geometric angles. The exact identities
\begin{equation}\label{occ:angle-identities}
 \begin{split}
 2d&=|n-u|^2+q^{-2},\\
 2D&=|n-a|^2+q_X^{-2},\\
 2e&=|a-u|^2+q^{-2}+q_X^{-2}
 \end{split}
\end{equation}
imply, on a nonempty bin,
\begin{equation}\label{occ:angle-ranges}
 \theta\geq1/p,\qquad \phi\geq c/w,\qquad
 \sigma\geq c/m,\qquad m=\min(p,w).
\end{equation}
All angles have at most $C\log(2+P)$ nonempty dyadic ranges.
For the relations between them, apply the triangle and reverse triangle
inequalities to
\[
 (u-n,q^{-1},0),\qquad(n-a,0,q_X^{-1})\in\R^5.
\]
Their norms are $\sqrt{2d},\sqrt{2D}$, and the norm of their sum is
$\sqrt{2e}$. Consequently $\sigma\leq C(\theta+\phi)$, and $\sigma$
is comparable to the larger of $\theta,\phi$ when those two angles are
sufficiently separated.

Fix a sufficiently small numerical $\kappa>0$.
The \emph{small sector} consists of representatives
$\theta<\kappa\phi$; here $\sigma\asymp\phi$, and we omit the
$\sigma$ partition entirely. All small-sector estimates below remain
valid on the enlarged range $\theta\leq16\kappa\phi$.
The \emph{remaining sector} has $\theta\geq\kappa\phi$
and retains the $\sigma$ partition. In particular it includes
$\theta\gg\phi$.
For the two sectors define
\begin{equation}\label{occ:sector-parameters}
 \begin{array}{c|c|c|c}
  &\nu&k_1&\omega\\ \hline
 \text{small}&\theta&1&\phi\\
 \text{remaining}&\min(\theta,\sigma)&1+\phi/\sigma&\sigma .
 \end{array}
\end{equation}
The parameter $\omega$ will be the velocity tolerance in the
occupation argument.

Let $N(t,s,n)$ be the integral of $f(s,X(t)-rn,v)$ over the momentum
restrictions of one bin, set to zero when the remaining restrictions
fail. Then
\begin{equation}\label{occ:density}
 N\leq C p^3\nu^2.
\end{equation}
Indeed the momentum shell has volume at most $Cp^3$. If $\theta$ is
small, \eqref{occ:angle-identities} restricts its velocity direction to a
spherical cap of radius $C\theta$ about $n$. If $\sigma$ is small,
both velocity magnitudes are close to one and their directions differ
by at most $C\sigma$, giving a cap about the receiver direction.
Use whichever cap is smaller. When the controlling parameter is bounded
below, the shell bound alone proves the same estimate with a larger
constant. The same argument shows that the allowed $n$ have spherical
area at most $C\phi^2$.

\subsection{Two changes of variables}

Define the occupation measure of the bin by
\begin{equation}\label{occ:measure}
 \mathsf M=\int dt\,ds\,dn\ r^2 D N(t,s,n),
\end{equation}
where $t$ ranges over the specified receiver times and $0\leq s<t$.
The constants below absorb $H_0$ in the comparison $r\asymp h$.
The factor $D$ makes this cone integral equal to particle mass integrated
over source-time hits, as the following changes of variables show.

\begin{lemma}[Crude occupation bounds]\label{occ:crude-lemma}
For every bin and every allowed set of receiver times,
\begin{align}
 \mathsf M&\leq C\frac{\max(I,h)}p,
                       \label{occ:energy}\\
 \mathsf M&\leq C I\frac{\phi^2}{\theta^2}.
                       \label{occ:time}
\end{align}
\end{lemma}

\begin{proof}
The characteristic change of variables below is used in
Pallard's work~\cite[Lemma~2.1]{Pallard2005}; see also
Luk--Strain~\cite[Proof of Proposition~5.2]{LukStrain2014Continuation}.
Fix $s$. The map
\[
 (t,n)\longmapsto y=X(t)-(t-s)n,\qquad t>s,
\]
is injective. Indeed, for fixed $y$, the function
$t-s-|X(t)-y|$ is strictly increasing: wherever differentiable its
derivative is at least $1-|a(t)|>0$. Its absolute Jacobian is $r^2D$.
This follows by taking two tangent vectors to the sphere and the time
derivative $a-n$, whose normal component is $-D$.
Hence \eqref{occ:measure} is the integral of $f(s,y,v)\,ds\,dy\,dv$
over the corresponding hits. The possible source times lie in an
interval of length at most $C(I+h)$. On the momentum bin $q\geq cp$,
so the particle energy bound proves \eqref{occ:energy}.

Passing at each fixed $s$ to the initial phase labels $z$, whose measure
is preserved by the characteristic flow, gives equivalently
\begin{equation}\label{occ:label-measure}
 \mathsf M=\int f_0(z)\,dz\int_{\text{hits of }z}ds.
\end{equation}
By Lemma~\ref{ret:receiver-time}, labels having a same-time collision
have zero $f_0(z)\,dz$ measure. Every remaining label has a single
increasing $C^1$ retarded branch, and
\begin{equation}\label{occ:retarded-derivative}
 \frac{dt}{ds}=\frac dD,\qquad \frac{ds}{dt}=\frac Dd.
\end{equation}
The branch begins when it enters the initial cone, if such an entry
occurs. Since its hit receiver times form a subset of $J$, their total
length is at most $I$. On the bin, $D/d\leq C\phi^2/\theta^2$.
Thus the hit duration of each label is at most
$CI\phi^2/\theta^2$. Integrating in \eqref{occ:label-measure} and using
the conserved finite particle mass proves \eqref{occ:time}.
\end{proof}

\paragraph{The three changes of variables.}
The preceding formulas use different fixed variables. For later reference,
their Jacobians are collected here:
\begin{center}
\renewcommand{\arraystretch}{1.4}
\begin{tabular}{@{}p{.15\linewidth}p{.47\linewidth}p{.30\linewidth}@{}}
\toprule
Fixed & Change of variables & Jacobian \\
\midrule
$(t,x)$ & Retarded $(y,v)$ to the initial label $z$
 & $dz=d\,dy\,dv$ \\
$s$ & $(t,n)\mapsto X(t)-(t-s)n$
 & $dy=r^2D\,dt\,dn$ \\
Label $z$ & $t=s+|X(t)-Y(s,z)|$
 & $dt/ds=d/D$ \\
\bottomrule
\end{tabular}
\end{center}
In particular the combined measure identity is
$ds\,f_0(z)dz=D\,dt\,f\,dy\,dv$.

\subsection{Stability on a short cell}

Fix large constants $C_0,C_1$, chosen below, and set
\begin{equation}\label{occ:lambda}
 \lambda=\left[
 C_1\frac Pm\left(\frac M\omega+
                         \sqrt{\frac{AL}{\omega}}\right)
                         \right]^{-2}.
\end{equation}
Call a bin \emph{actual} if $\omega>C_0/m$, and \emph{deficit}
otherwise. In an actual bin \eqref{occ:angle-identities} implies
$|u(s)-a(t)|\geq c\omega$ at each hit: the two inverse-momentum
terms are smaller than a fixed small multiple of $\omega^2$ once
$C_0$ is large. In the small sector use additionally
$\theta\leq16\kappa\phi$.
In a deficit bin,
\begin{equation}\label{occ:deficit-lambda}
 \lambda\leq\frac{C_0^2}{C_1^2P^2}.
\end{equation}

We record explicitly the stability supplied by the bootstrap. Given a
fixed constant $C_*$ and a hit, all source times within $C_*\lambda$
of its source time, and all receiver times within $C_*\lambda$ of its
receiver time, clipped to $[0,\tau]$, satisfy
\begin{equation}\label{occ:stability}
 |\Delta V|\leq\varepsilon_0 m\omega,
 \qquad |\Delta u|\leq\varepsilon_0\omega,
 \qquad q\asymp q_{\rm hit}.
\end{equation}
Here $\varepsilon_0>0$ can be prescribed by choosing $C_1$ sufficiently
large in terms of $C_*$ and the fixed comparisons, independently of
$M,A$.
To prove this, choose a dyad $w'$ nearest $q_{\rm hit}$ and stop at the
first exit from $[q_{\rm hit}/2,2q_{\rm hit}]$. Before exit the bootstrap
applies with $w'\asymp q_{\rm hit}\geq cm$. Its two terms on an interval
of length at most $C_*\lambda$ are bounded by
\[
 MP\sqrt{C_*\lambda}\leq
       \frac{\sqrt{C_*}}{C_1}m\omega,
 \qquad
 C A C_*\lambda\frac{P^2L}{m}
       \leq\frac{CC_*}{C_1^2}m\omega.
\]
For sufficiently large $C_1$, this excludes the first exit, because
$q$ is 1-Lipschitz and $\omega\leq1$. The same stopping argument runs
backwards from the hit. The segment between the two momentum vectors
has $q\geq cq_{\rm hit}$; integrating $\|D_vu\|\leq1/q$ on that segment
then proves the velocity estimate. This proves \eqref{occ:stability}.
We now fix $C_*$ large enough to cover the cell and retarded-time
separations used below, and then fix $C_1$ accordingly.

\begin{lemma}[Occupation on stable cells]\label{occ:cell-lemma}
In every actual bin,
\begin{equation}\label{occ:actual}
 \mathsf M\leq C\frac{hk_1}{p}
                       \max(1,I/\lambda,h/\lambda).
\end{equation}
\end{lemma}

\begin{proof}
Suppose first $h\leq\lambda$, and partition $[0,\tau]$ into intervals
of length $\lambda$, with the last interval clipped if necessary.
For upper bounds we may enlarge the receiver times to all eligible times
in $J$ and use closed bin ranges. The labels with some hit in a closed
cell form a measurable set: it is the projection of a compact subset of
the product of the compact initial support, the cell, $J$, and $S^2$;
the fixed bin has $r\geq ch>0$.

For any label with a hit $(s_0,t_0)$ in the cell, use that hit as a
reference. At every same-time point $s$ of the cell, both $u(s,z)$ and
$a(s)$ are within a small multiple of $\omega$ of their respective
reference velocities: the time separations are at most
$\lambda+C_{T_0}h\leq C_*\lambda$.
The actual velocity difference therefore has projection at least
$c\omega$ in one fixed direction throughout the cell. This direction
may depend on the label and its reference hit.
At any hit in the cell,
\[
 Y(s,z)-X(s)=\int_s^t(a(t')-n)\,dt',
 \qquad |Y(s,z)-X(s)|\leq Ch(\phi+\omega),
\]
using $|a(t)-n|\leq C\phi$ and receiver stability between $s$ and $t$.
The projected same-time separation is strictly monotone with derivative
of magnitude at least $c\omega$. Its values at hits lie in an interval
of length $Ch(\phi+\omega)$, so their total $s$-duration is at most
$Ch(1+\phi/\omega)\leq Chk_1$. This holds even if the hit set is
disconnected.

All labels that hit the cell have $q\asymp p$ at any one common time in
the cell, by stability. Conservation of phase volume and the energy
bound give their total label mass at most $C/p$. The source-time range
of all hits has length at most $C(I+h)$, so at most
$C\max(1,I/\lambda,h/\lambda)$ cells can occur. Formula
\eqref{occ:label-measure} proves \eqref{occ:actual}.
If $h>\lambda$, instead use \eqref{occ:energy}: its right-hand side is
bounded by that of \eqref{occ:actual}, since $k_1\geq1$ and
$h\max(I,h)/\lambda\geq\max(I,h)$.
\end{proof}

\subsection{One energy budget for all distance and source-angle bins}

Choose a fixed enlargement
\[
 J^*=\{s\in[0,\tau]:\operatorname{dist}(s,J)\leq C_{T_0}I\},
\]
with its constant large enough for every source cell considered below
when $h\leq\lambda\leq I$. In particular $|J^*|\leq CI$.
For fixed $s,y$, let $t_{\rm ret}(s,y)\in(s,\tau]$ be the unique solution
of $t-s=|X(t)-y|$, if it exists; omit its tip $y=X(s)$.
This definition deliberately imposes no restriction that
$t_{\rm ret}$ belong to $J$ or to the eligible receiver subset.

The first indicator below is used to average the mass of hit labels over
a stable cell, using the velocity separation that persists there.
The second records the receiver angle at a retarded hit and is used
when the spatial change of variables estimates hits directly.
For each $p,\phi$ define these indicators with fixed sufficiently enlarged
comparison constants:
\[
 \begin{split}
 H^{(1)}_{p\phi}(s,y,v)
  &=\mathbf1\{q(v)\asymp p, |u(v)-a(s)|\asymp\phi\},\\
 H^{(2)}_{p\phi}(s,y,v)
  &=\mathbf1\left\{q(v)\asymp p, t_{\rm ret}\text{ exists},
    \sqrt{\frac{1-n_{\rm ret}\cdot a(t_{\rm ret})}{2}}\asymp\phi
                    \right\},\\
 n_{\rm ret}&=\frac{X(t_{\rm ret})-y}{t_{\rm ret}-s}.
 \end{split}
\]
The second indicator is zero when its retarded time does not exist.
The strictly monotone implicit equation makes its non-tip solution
continuous when the root lies in $(s,\tau)$ and continuously extendible
to a root at $\tau$. The existence condition is a closed inequality
at $\tau$, after excluding the tip; thus the indicators are measurable.
Set
\begin{equation}\label{occ:energy-budget}
 \mathcal E_{p\phi}=\frac1I\int_{J^*}\int_{\R^3\times\R^3}
 qf\bigl(H^{(1)}_{p\phi}+H^{(2)}_{p\phi}\bigr)\,dy\,dv\,ds.
\end{equation}
These numbers depend on the receiver, $J,I,p,\phi$, but not on
$\theta,h$ or on the eligible measurable subset. Their comparison
constants, and hence the bound
\begin{equation}\label{occ:energy-budget-sum}
 \sum_{p,\phi}\mathcal E_{p\phi}\leq C,
\end{equation}
are independent of $M,A$. Indeed at each $(s,y,v)$ only a bounded number
of $p$ and $\phi$ dyads can occur in either family. For the second family
there is just one retarded time. Summing the indicators and using the
energy bound and $|J^*|\leq CI$ proves \eqref{occ:energy-budget-sum}.

\begin{proposition}[Baseline occupation with a joint energy budget]\label{occ:refined-proposition}
Suppose $h\leq I$, excluding only deficit bins with $h\leq\lambda$.
Then
\begin{equation}\label{occ:refined}
 \mathsf M\leq C\frac{hk_1}{p}
             \left(1+\frac I\lambda\mathcal E\right),
 \qquad
 \mathcal E=\begin{cases}
  \mathcal E_{p\phi},&\text{small sector},\\
  1,&\text{remaining sector}.
 \end{cases}
\end{equation}
Consequently, with
$R=1+\sqrt{I/\lambda}\sqrt{\mathcal E}$,
\begin{equation}\label{occ:R-bound}
 \mathsf M\leq C\frac{hk_1}{p}R^2
\end{equation}
for all $h$, except deficit bins with $h\leq\min(I,\lambda)$.
\end{proposition}

\begin{proof}
In the remaining sector, \eqref{occ:actual} proves the claim for actual
bins; if $h>\lambda$, \eqref{occ:energy} gives
$\mathsf M\leq CI/p\leq Ch(I/\lambda)/p$. This also treats the allowed
deficit bins.

Consider the small sector. When $h\leq\lambda\leq I$ and the bin is
actual, denote the set of hit labels in a cell $C$ by $Z_C$. At every
time $s\in C$ all these labels have both $q\asymp p$ and
$|u-a(s)|\asymp\phi$, by the same stability argument as in
Lemma~\ref{occ:cell-lemma}. For a full cell, averaging its mass over its
whole length therefore gives
\[
 \int_{Z_C}f_0(z)\,dz
 \leq\frac C{p\lambda}\int_C\int qf
                  H^{(1)}_{p\phi}\,dy\,dv\,ds.
\]
All such cells lie in $J^*$ and their interiors are disjoint. The at most
two cells clipped by the endpoints of $[0,\tau]$ have mass at most $C/p$
by the common-time argument, without averaging. Each label has hit
duration at most $Ch$ in its cell. Summing
\eqref{occ:label-measure} proves
$\mathsf M\leq Chp^{-1}(1+(I/\lambda)\mathcal E_{p\phi})$.
If $\lambda>I$, the source-time range has length at most $CI$ and meets
only a bounded number of cells. Their crude mass and duration bounds
already give $\mathsf M\leq Ch/p$; no averaging outside $J^*$ is used.

Finally suppose $\lambda<h\leq I$, whether actual or deficit. In the
space map of Lemma~\ref{occ:crude-lemma}, each hit belongs to
$H^{(2)}_{p\phi}$ and its source time belongs to $J^*$. Hence
\[
 \mathsf M\leq\frac C p\int_{J^*}\int qf
                    H^{(2)}_{p\phi}\,dy\,dv\,ds
       \leq\frac{CI}p\mathcal E_{p\phi}.
\]
Since $h/\lambda>1$, this is bounded by the right-hand side of
\eqref{occ:refined}.
The estimate $1+(I/\lambda)\mathcal E\leq R^2$ proves
\eqref{occ:R-bound} when $h\leq I$. For $h>I$,
\eqref{occ:energy} gives $\mathsf M\leq Ch/p$, and
$k_1R^2\geq1$ completes the proof. The budget
\eqref{occ:energy-budget} can be defined for every bin, regardless of
whether $h\leq I$; its use in this last case requires only nonnegativity.
\end{proof}
